%!TEX root = ../main_thesis.tex

\chapter{Esempi di modelli}
\label{app:esempi}

A titolo di esempio completo si riporta il modello del \emph{problema della
dieta}, che impiega gran parte dei costrutti del linguaggio descritti nel
Capitolo~\ref{cap:linguaggio}: costanti scalari e array (anche
multidimensionali), funzioni della libreria standard (\texttt{enumerate},
\texttt{len}), sommatorie indicizzate, variabili composte, vincoli nominati e
replicati tramite la clausola \texttt{for}, e variabili con estremi dipendenti
dai dati. L'obiettivo è minimizzare il costo complessivo della dieta mantenendo
ciascun nutriente entro i limiti previsti.

\begin{lstlisting}[language=rooc, caption={Il problema della dieta in ROOC.}, label={lst:diet}]
min sum((cost, i) in enumerate(C)) { cost * x_i }
s.t.

    min_{nutrient[j]}: //diet must have at least of nutrient j
        sum(i in 0..f) { a[i][j] * x_i} >= Nmin[j] for j in 0..len(Nmin)

    max_{nutrient[j]}: //diet must have at most of nutrient j
        sum(i in 0..f) { a[i][j] * x_i } <= Nmax[j] for j in 0..len(Nmax)

where
    // Cost of chicken, rice, avocado
    let C = [1.5, 0.5, 2.0]
    // Min and max of: protein, carbs, fats
    let nutrient = ["protein", "carbs", "fats"]
    let Nmin = [50, 200, 0]
    let Nmax = [150, 300, 70]
    // Min and max servings of each food
    let Fmin = [1, 1, 1]
    let Fmax = [5, 5, 5]
    let a = [
        //protein, carbs, fats
        [30, 0, 5], // Chicken
        [2, 45, 0], // Rice
        [2, 15, 20] // Avocado
    ]
    // Number of foods
    let f = len(a)
    // Number of nutrients
    let n = len(Nmax)
define
    //bound the amount of each serving of food i
    x_i as NonNegativeReal(Fmin[i], Fmax[i]) for i in 0..n
\end{lstlisting}
